\documentclass[11pt,a4paper]{article}
\usepackage[T1]{fontenc}
\usepackage[utf8]{inputenc}
\usepackage{lmodern,amsmath,amssymb}
\usepackage[margin=25mm]{geometry}
\usepackage[hidelinks]{hyperref}
\hypersetup{pdftitle={The record costs recoil: a floor for every probe state},pdfauthor={navstokgap research working notes}}
\newcommand{\tightlist}{\setlength{\itemsep}{0pt}\setlength{\parskip}{0pt}}
\setlength{\emergencystretch}{3em}
\setcounter{secnumdepth}{0}
\begin{document}
\hypertarget{the-record-costs-recoil-a-floor-for-every-probe-state}{%
\section{The record costs recoil: a floor for every probe
state}\label{the-record-costs-recoil-a-floor-for-every-probe-state}}

For every protocol of momentum-transfer marks, with probes in any states
whatever (Gaussian or not, squeezed, correlated, mixed, grid states),
deciding between free motion and a constant force \(F\) over a duration
\(\tau\) at error probability \(\epsilon\) requires

\[\boxed{\;s\sum_j\Delta_j\ \ge\ 8\,(1-2\epsilon)\,\hbar,\qquad s=\frac{F\tau^2}{2m},\;}\]

where \(s\) is Newton's sagitta for the interval and \(\Delta_j\) is the
spread of the impulse delivered by mark \(j\). The resolutions of the
marks do not enter. For a general force the sagitta is replaced by the
uniform distance of the forced displacement from straight motion
(Theorem R).

The theorem answers the question the non-Gaussian probes raise. The
floor \(\tau\Delta E\ge24z^2\hbar\sqrt{(1-\rho)/(1+\rho)}\) of the
\href{mark-cost-and-statistical-floor.md}{mark-cost note} is a theorem
about Gaussian probes: a grid state with zero position--momentum
correlation decides the comparison at any force (Proposition G), so no
bound on second moments alone survives. What does survive, for every
probe, is the recoil. A pointer registers a displacement only through
the spread of its conjugate momentum, and that momentum is the impulse
the body receives. This is M3 in its universal form: a record of the
sagitta must leave the delivered impulse undetermined by at least
\(8(1-2\epsilon)\hbar/s\) in total. The correlated three-mark protocol
of Proposition D attains the bound within a factor
\(\sqrt2\,z_{1-\epsilon}/(1-2\epsilon)\), uniformly in its squeezing
(Section 4). Exploratory; no ledger promotion.

\hypertarget{setting}{%
\subsection{1. Setting}\label{setting}}

Transverse coordinate \(\hat y\), momentum \(\hat p\), mass \(m\), on
\([0,\tau]\). Marks \(j=1,\dots,k\) act at times \(t_j\): the body is
coupled to probe \(j\), prepared in a state \(\varphi_j\) (pure or
mixed, arbitrary), by
\(U_j=\exp(-\frac i\hbar\lambda_j\hat y\hat P_j)\), and the pointer
\(\hat Q_j\) is read afterwards with no further evolution. Under
hypothesis \(\mathrm I\) the body moves freely; under \(\mathrm F\) it
feels a force whose displacement relative to free motion from rest at
the origin is \(P(t)\), with \(P(0)=P'(0)=0\); for the constant force
\(P(t)=Ft^2/2m\). The body's initial state is unknown: a test
\textbf{decides at error \(\epsilon\)} when its equal-prior error
probability is at most \(\epsilon\) for every initial state under
\(\mathrm I\) and every initial state under \(\mathrm F\). Write
\(\hat N_j=\hat Q_j/\lambda_j\) and \(\hat D_j=-\lambda_j\hat P_j\) for
the error and the delivered impulse of mark \(j\), as in Theorem A, and
\(\Delta_j=\lambda_j\Delta_{\varphi_j}\hat P_j\) for its recoil spread.

\textbf{Lemma 1 (the record).} In the Heisenberg picture the \(j\)-th
reading is

\[\hat R_j=\hat y+\frac{\hat p\,t_j}{m}+\theta P(t_j)
+\sum_{i:\,t_i<t_j}\frac{t_j-t_i}{m}\hat D_i+\hat N_j ,\]

with \(\theta=0\) under \(\mathrm I\) and \(\theta=1\) under
\(\mathrm F\), and the initial pointer \(\hat Q_j\) appears in
\(\hat R_j\) alone.

\emph{Proof.} Theorem A of the mark-cost note for each coupling, the
linear free evolution between marks, and the c-number force, which adds
\(P(t)\) to \(\hat y(t)\). The pointer \(\hat Q_j\) is shifted by the
coupling and coupled to nothing afterwards. \(\square\)

\textbf{Lemma 2 (the signal moves into the pointers).} For real \(a,b\)
let the body's initial state under \(\mathrm F\) be the state under
\(\mathrm I\) translated by \(-a\) in position and \(-mb\) in momentum.
Then the record under \(\mathrm F\) has the distribution of the record
under \(\mathrm I\) with each probe state \(\varphi_j\) replaced by its
translate by \(\lambda_jc_j\) in \(\hat Q_j\), where
\(c_j=P(t_j)-a-bt_j\).

\emph{Proof.} The translation of the body replaces \(\hat y,\hat p\) in
Lemma 1 by \(\hat y-a\), \(\hat p-mb\), which shifts \(\hat R_j\) by
\(-a-bt_j\); the force adds \(P(t_j)\). The net effect is the c-number
shift \(\hat R_j\mapsto\hat R_j+c_j\). Translating probe \(j\) by
\(\lambda_jc_j\) in \(\hat Q_j\) shifts \(\hat N_j\) by \(c_j\) and
leaves every \(\hat D_i\) and every other reading unchanged, by the last
clause of Lemma 1. The record operators commute, so their joint
distribution is fixed by these shifts and the states. \(\square\)

\hypertarget{the-theorem}{%
\subsection{2. The theorem}\label{the-theorem}}

\textbf{Theorem R.} Every protocol deciding at error \(\epsilon\)
satisfies

\[1-2\epsilon\ \le\ \frac1\hbar\,\min_{a,b}\sum_j\Delta_j\,\bigl|P(t_j)-a-bt_j\bigr| .\]

\emph{Proof.} Fix \(a,b\) and any initial state \(\sigma\) under
\(\mathrm I\), and take under \(\mathrm F\) the translate of Lemma 2. A
test with error at most \(\epsilon\) against this pair has
\(\mathrm{TV}\ge1-2\epsilon\) between the two record distributions. By
Lemma 2 those are the records of one fixed protocol, with body state
\(\sigma\), applied to the probe states \(\otimes_j\varphi_j\) and
\(\otimes_j\varphi_j'\), \(\varphi_j'\) the translate by
\(q_j=\lambda_jc_j\). The data-processing inequality and the triangle
inequality over the factors give

\[\mathrm{TV}\le\tfrac12\bigl\|\otimes_j\varphi_j-\otimes_j\varphi_j'\bigr\|_1
\le\sum_j\tfrac12\|\varphi_j-\varphi_j'\|_1 .\]

The translate is
\(e^{-iq_j\hat P_j/\hbar}\varphi_je^{iq_j\hat P_j/\hbar}\). By
Mandelstam--Tamm for a pure state, and through a purification for a
mixed one exactly as in Theorem 2 of the
\href{planck-gap-probabilistic.md}{probabilistic note},
\(\frac12\|\varphi_j-\varphi_j'\|_1\le|q_j|\Delta_{\varphi_j}\hat P_j/\hbar=\Delta_j|c_j|/\hbar\).
Minimize over \(a,b\). \(\square\)

\textbf{Corollary R1 (the Galileo comparison).} For \(P(t)=Ft^2/2m\),
the line \(a+bt\) that best approximates \(P\) uniformly on \([0,\tau]\)
is \(\frac{F}{2m}(\tau t-\tau^2/8)\), and
\(|P(t)-a-bt|\le F\tau^2/16m=s/8\) there, with equality at
\(t=0,\tau/2,\tau\). Hence

\[s\sum_j\Delta_j\ \ge\ 8(1-2\epsilon)\hbar .\]

For a general force the same step gives
\(E_1(P)\sum_j\Delta_j\ge(1-2\epsilon)\hbar\), with
\(E_1(P)=\min_{a,b}\max_{[0,\tau]}|P(t)-a-bt|\) the uniform distance of
the forced displacement from straight motion. \(\square\)

The mechanism is the conjugacy of Theorem A read in the other direction.
The record of the force is a translation of the pointers, and a pointer
feels a translation only through the spread of the operator that
generates it, which is its momentum \(\hat P_j\). That same momentum,
times \(\lambda_j\), is the impulse the body receives. So sensitivity is
paid in recoil at the exchange rate \(\hbar\), whatever the pointer's
state, and a resolution bought by squeezing, correlation or a grid is
bought with recoil.

For adaptive protocols, in which couplings and times depend on earlier
readings, the hybrid argument of the probabilistic note gives the same
bound with each \(\Delta_j|c_j|\) replaced by its supremum over the
choices available at step \(j\).

\hypertarget{the-gaussian-floor-needs-gaussian-probes}{%
\subsection{3. The Gaussian floor needs Gaussian
probes}\label{the-gaussian-floor-needs-gaussian-probes}}

\textbf{Proposition G.} In the three-mark protocol of Proposition D of
the mark-cost note (marks at \(0,\tau/2,\tau\), test \(R_1-2R_2+R_3\),
outer marks sharp), let the middle probe be a grid state with a real
wavefunction, so that \(\operatorname{Cov}(\hat Q_2,\hat P_2)=0\) and
its correlation \(\rho\) vanishes. For every \(F\) whose signal
\(u^{\mathsf T}P=F\tau^2/4m\) is not a multiple of
\(\sigma=\sqrt{2\pi\hbar\,|u_2S_2|/m}=\sqrt{2\pi\hbar\tau/m}\), the
error probability tends to zero as the grid is made finer.

\emph{Proof.} The statistic's noise is
\(-2\hat N_2+\frac{\tau}{2m}\hat D_2\) plus the vanishing contributions
of the outer marks, that is \(\alpha\hat Q_2+\beta\hat P_2\) with
\(\alpha\beta=|u_2S_2|/m=\tau/m\) independent of \(\lambda_2\). Take the
ideal grid state stabilized by the commuting translations
\(e^{-i\ell_Q\hat P_2/\hbar}\) and \(e^{i\ell_P\hat Q_2/\hbar}\),
\(\ell_Q\ell_P=2\pi\hbar\) (Gottesman, Kitaev and Preskill,
\href{https://doi.org/10.1103/PhysRevA.64.012310}{PRA \textbf{64},
012310, 2001}, metadata), with the aspect chosen as
\(\ell_P/\ell_Q=\alpha/\beta\). The product of the two stabilizers is
\(-e^{ic(\alpha\hat Q_2+\beta\hat P_2)}\) with
\(c=\ell_P/(\hbar\alpha)\), by the Baker--Campbell--Hausdorff phase
\(e^{-i\pi}\), so \(\alpha\hat Q_2+\beta\hat P_2\) is supported on a
comb of spacing \(2\pi/c=\sqrt{2\pi\hbar\alpha\beta}=\sigma\).
Finite-energy grid states with Gaussian teeth of width \(w\) and a
Gaussian envelope keep a real wavefunction and converge to the ideal
state as \(w\to0\); the statistic then has teeth of width of order
\(w\), and a shift by a non-multiple of \(\sigma\) separates the two
combs. \(\square\)

Such states measure both quadratures of a small displacement at once,
which is the principle of the grid-state displacement sensor of
Duivenvoorden, Terhal and Weigand
(\href{https://doi.org/10.1103/PhysRevA.95.012305}{PRA \textbf{95},
012305, 2017}, metadata). Two features connect it to the rest of the
programme. Its recoil diverges: the envelope of the momentum
distribution widens as \(\hbar/w\), so Theorem R is met. And its blind
spots are the shifts in \(\sigma\mathbb Z\), the translations whose Weyl
phase with the stabilizers is a multiple of \(2\pi\): the grid reads the
signal modulo the same central phase that the
\href{polygon-lift-phase.md}{polygon-lift note} identifies as the part
no preparation can move.

So Theorems B and C of the mark-cost note hold for Gaussian probes,
where the Wigner function is a density and the record is the classical
Gaussian model, and the correlation coefficient is the resource that
lowers the floor within that class. Outside it the second moments do not
bound the error, and Theorem R is the statement that holds.

\hypertarget{tightness}{%
\subsection{4. Tightness}\label{tightness}}

In the correlated protocol of Proposition D the outer marks are sharp,
so their recoils are unbounded and the minimum in Theorem R takes the
line through the endpoints, \(c_1=c_3=0\), \(|c_2|=F\tau^2/8m\). The
balanced middle mark has \(2\delta_2=\tau\Delta_2/2m\) and, for a pure
Gaussian probe, \(\delta_2\Delta_2=\hbar/(2\sqrt{1-\rho^2})\), so
\(\Delta_2^2=2m\hbar/(\tau\sqrt{1-\rho^2})\). At its decision threshold
\(\tau\Delta E=32z^2\hbar\sqrt{(1-\rho)/(1+\rho)}\),

\[\Delta_2|c_2|=z_{1-\epsilon}\,\hbar\,\sqrt{\frac{2}{1+\rho}},\]

between \(z\hbar\) and \(\sqrt2z\hbar\) for every \(\rho\in[0,1)\).
Theorem R asks for \((1-2\epsilon)\hbar\), so the protocol sits within
the factor \(\sqrt2z_{1-\epsilon}/(1-2\epsilon)\) of the bound, about
\(2.6\) at five per cent error, however strongly it is squeezed. The
recoil-weighted sagitta is the right universal resource: squeezing moves
the protocol along it and never below it.

\hypertarget{what-this-changes}{%
\subsection{5. What this changes}\label{what-this-changes}}

\textbf{The three floors, ordered.} Resource-free: the phase
\(\mathcal K_\tau/\hbar\) between the motion and its inscribed polygon
(polygon-lift note, Theorems 1 and 7). Universal over probe states: the
recoil bound \(s\sum_j\Delta_j\ge8(1-2\epsilon)\hbar\) (Theorem R).
Universal over instruments with a bounded laboratory: the aperture bound
\(F\tau L+\frac{F\tau^2}{2m}P\ge(1-2\epsilon)\hbar\) (probabilistic
note). And for Gaussian probes with correlation at most \(\rho\), the
sharp floor \(\tau\Delta E\ge24z^2\hbar\sqrt{(1-\rho)/(1+\rho)}\)
(mark-cost note).

\textbf{M3, restated.} The mark-floor note's premise has two clauses, a
product of resolution and recoil and the independence of the two.
Theorem R shows that the recoil clause alone is universal: a mark that
records the sagitta must leave the delivered impulse undetermined, with
\(s\sum\Delta_j\ge8(1-2\epsilon)\hbar\). In Newton's words from the
\href{newton-mark-floor.md}{mark-floor note}, the premise that carries
\(h>0\) for every probe is that a record costs an undetermined impulse,
and the pairing is with the sagitta of Lemma X.

\textbf{Beyond the momentum-transfer class.} The proof uses only that
the pointer's error is translated by the unitary generated by the
delivered impulse, \([\hat N_j,\hat D_j]=-i\hbar\), and that the
apparatus enters the record through \(\hat N_j\) and \(\hat D_j\). The
commutation of the pointer with the body's momentum after any coupling
gives \([\hat N,\hat D]=-i\hbar-[\hat y,\hat D]-[\hat N,\hat p]\), so
Theorem R extends to every instrument whose error and recoil do not
depend on the body (Arthurs and Goodman,
\href{https://doi.org/10.1103/PhysRevLett.60.2447}{PRL \textbf{60},
2447, 1988}, metadata). That is the next step, and the body-dependent
case, where Ozawa's relation
(\href{https://doi.org/10.1103/PhysRevA.67.042105}{PRA \textbf{67},
042105, 2003}, metadata) adds the body's own spreads, is the one after.

\hypertarget{consequence-for-state}{%
\subsection{6. Consequence for STATE}\label{consequence-for-state}}

Theorem A's extension, option 3 of the 2026-09-22 plan, is settled
constructively. The Gaussian floor holds for Gaussian probes and grid
probes break it (Proposition G); the universal statement over probe
states is Theorem R, sagitta times total recoil at least
\(8(1-2\epsilon)\hbar\), attained within a constant by squeezed probes
at every squeezing. Next: option 1 (additive-noise instruments), then
option 2 (Ozawa's relation).
\end{document}
